\documentclass[11pt,a4paper]{article}
\usepackage[utf8]{inputenc}
\usepackage[T1]{fontenc}
\usepackage[spanish,es-minimal,es-noshorthands]{babel}
\usepackage{lmodern}
\linespread{1.22}
\usepackage[a4paper,top=2.4cm,bottom=2.2cm,left=1.5cm,right=1.5cm,headheight=1.1cm,headsep=0.5cm,footskip=1cm]{geometry}
\usepackage{amsmath,amssymb,cancel}
\usepackage{icomma}
\usepackage{graphicx}
\usepackage[table]{xcolor}
\usepackage{tikz}
\usepackage{pgfplots}\pgfplotsset{compat=1.17}
\usepackage{fontawesome5}
\usepackage{ccicons}
\usepackage{fancyhdr}
\usepackage{enumitem}
\usepackage{booktabs}
\usepackage{tabularx}
\usepackage{multirow}
\usepackage{needspace}
\usepackage[most]{tcolorbox}
\usepackage{titlesec}
\usepackage[hidelinks]{hyperref}
\hypersetup{pdftitle={Apuntes de teoría. Unidad 10: Distribución binomial y normal. 2.º de Bachillerato CCSS},pdfsubject={Matemáticas Aplicadas a las Ciencias Sociales II},pdfkeywords={variable aleatoria, función de probabilidad, esperanza, distribución binomial, número combinatorio, distribución normal, campana de Gauss, normal tipificada, tipificación, aproximación de la binomial por la normal, corrección de continuidad}}

% colors de la web (FentMates, bloc de probabilitat: morat; es manté el nom fmgreen dels apunts anteriors)
\definecolor{fmnavy}{HTML}{1B2748}
\definecolor{fmgreen}{HTML}{684E88}
\definecolor{fmgreenl}{HTML}{8064A1}
\definecolor{fmyellow}{HTML}{FFC94D}
\definecolor{fmgold}{HTML}{E9A919}
\definecolor{fmaccent}{HTML}{FF7A50}
\definecolor{fmteal}{HTML}{2BB3A3}
\definecolor{fmdef}{HTML}{1D4F73}\definecolor{fmdefbg}{HTML}{EEF4F9}
\definecolor{fmprop}{HTML}{3B4252}\definecolor{fmpropbg}{HTML}{F3F4F7}
\definecolor{fmex}{HTML}{A3203A}\definecolor{fmexbg}{HTML}{FBF0F2}
\definecolor{fmmuted}{HTML}{5D6685}
\definecolor{fmsurface}{HTML}{EEF1F8}

\newcommand{\lt}{<}\newcommand{\gt}{>}
% fórmula en línia que no es parteix mai; si no cap en una línia, es posa centrada
\newcommand{\fmm}[1]{\setbox0=\hbox{$#1$}\ifdim\wd0>\linewidth\[#1\]\else\mbox{\box0}\fi}
\DeclareUnicodeCharacter{2212}{-}
\setlength{\parindent}{0pt}
\setlength{\parskip}{12pt}
\setlength{\lineskip}{5pt}\setlength{\lineskiplimit}{5pt}  % més aire al voltant de les fraccions en línia
\setlength{\emergencystretch}{3em}
\setlist{itemsep=4pt,topsep=4pt}
\renewcommand{\thesection}{10.\arabic{section}}
\titleformat{\section}{\color{fmnavy}\Large\bfseries}{\colorbox{fmdef}{\color{white}\strut\ \thesection\ }}{0.6em}{}[{\color{fmdef}\titlerule[1.2pt]}]
\titlespacing*{\section}{0pt}{16pt}{8pt}
\titleformat*{\subsection}{\color{fmgreen}\large\bfseries}
\allowdisplaybreaks

% logo de la web: adhesiu verd amb la gràfica + «Fent» «Mates»
\newcommand{\fmlogo}[1]{\begin{tikzpicture}[x=#1/124,y=#1/124,baseline=-0.35*#1]
  \fill[fmgreenl,rounded corners=#1*0.24] (-62,-62) rectangle (62,62);
  \begin{scope}[yscale=-1,scale=1.15,line width=#1*0.045,line cap=round,line join=round,draw=white]
  \draw (-36,24)--(36,24) (-28,32)--(-28,-32);
  \draw (-34,18) .. controls (-14,18) and (-10,-24) .. (6,-24) .. controls (22,-24) and (24,10) .. (36,-6);
  \end{scope}\end{tikzpicture}}
% dibuix de la targeta de la unitat a la web (illustrations.js, «normal»), en blanc
\newcommand{\fmcard}{\begin{tikzpicture}[x=0.058cm,y=-0.058cm,white,line width=4pt,line cap=round,line join=round]
  \draw[line width=3pt,opacity=0.7] (10,104)--(230,104);
  \begin{scope}\clip (14,104)--(14,103) .. controls (50,101) and (70,92) .. (88,64) .. controls (102,40) and (112,22) .. (120,22) .. controls (128,22) and (138,40) .. (152,64) .. controls (170,92) and (190,101) .. (226,103)--(226,104)--cycle;
    \fill[fill opacity=0.2] (134,0) rectangle (196,104);\end{scope}
  \draw (14,103) .. controls (50,101) and (70,92) .. (88,64) .. controls (102,40) and (112,22) .. (120,22) .. controls (128,22) and (138,40) .. (152,64) .. controls (170,92) and (190,101) .. (226,103);
  \draw[line width=3pt,dash pattern=on 6pt off 7pt] (120,22)--(120,104) (134,34)--(134,104);
  \fill[fill opacity=0.35] (20,88) rectangle (30,104) (32,74) rectangle (42,104) (44,66) rectangle (54,104) (56,78) rectangle (66,104);
  \node[anchor=base west,inner sep=0pt,font=\fontsize{22}{22}\selectfont\itshape] at (101,14) {N(\(\mu\),\(\sigma\))};
  \node[anchor=base west,inner sep=0pt,font=\fontsize{19}{19}\selectfont\itshape] at (20,60) {B(n,p)};
\end{tikzpicture}}
\newcommand{\fmbrand}{\raisebox{-0.25\height}{\includegraphics[height=0.5cm]{logo-fentmates.pdf}}}

% capçalera i peu
\pagestyle{fancy}\fancyhf{}
\fancyhead[L]{\small\fmbrand}
\fancyhead[R]{\small\color{fmmuted}2.º Bachillerato CCSS · Unidad 10: Distribución binomial y normal}
\fancyfoot[L]{\scriptsize\color{fmmuted}\href{https://creativecommons.org/licenses/by-nc-sa/4.0/deed.ca}{\ccbyncsa}\ \ FentMates · fentmates.org · CC BY-NC-SA 4.0}
\fancyfoot[R]{\small\color{fmmuted}\thepage}
\renewcommand{\headrulewidth}{0.4pt}\renewcommand{\headrule}{\color{fmmuted!40}\hrule height \headrulewidth}
\fancypagestyle{primera}{\fancyhf{}\renewcommand{\headrulewidth}{0pt}
  \fancyfoot[L]{\scriptsize\color{fmmuted}\href{https://creativecommons.org/licenses/by-nc-sa/4.0/deed.ca}{\ccbyncsa}\ \ FentMates · fentmates.org · CC BY-NC-SA 4.0}
  \fancyfoot[R]{\small\color{fmmuted}\thepage}}

% caixes (mateixos colors que la web)
% quan una caixa es parteix entre dues pàgines, la vora del tall és discontínua
\tcbset{fmbase/.style={enhanced,breakable,arc=3mm,boxrule=1pt,left=3.5mm,right=3.5mm,top=2mm,bottom=2.5mm,before skip=10pt,after skip=10pt,fonttitle=\bfseries,
  frame style={fill=tcbcolframe},
  overlay first={\fill[tcbcolback] ([xshift=1pt]frame.south west) rectangle ([xshift=-1pt,yshift=1.2pt]frame.south east);
    \draw[tcbcolframe,line width=1pt,dash pattern=on 5pt off 3pt] (frame.south west)--(frame.south east);},
  overlay middle={\fill[tcbcolback] ([xshift=1pt]frame.south west) rectangle ([xshift=-1pt,yshift=1.2pt]frame.south east);
    \fill[tcbcolback] ([xshift=1pt]frame.north west) rectangle ([xshift=-1pt,yshift=-1.2pt]frame.north east);
    \draw[tcbcolframe,line width=1pt,dash pattern=on 5pt off 3pt] (frame.south west)--(frame.south east) (frame.north west)--(frame.north east);},
  overlay last={\fill[tcbcolback] ([xshift=1pt]frame.north west) rectangle ([xshift=-1pt,yshift=-1.2pt]frame.north east);
    \draw[tcbcolframe,line width=1pt,dash pattern=on 5pt off 3pt] (frame.north west)--(frame.north east);}}}
\newtcolorbox{fmdef}[1]{fmbase,colback=fmdefbg,colframe=fmdef,coltitle=white,title={#1}}
\newtcolorbox{fmprop}[1]{fmbase,colback=fmpropbg,colframe=fmprop,coltitle=white,title={#1},halign title=center}
\newtcolorbox{fmex}[1]{fmbase,before upper=\raggedright,colback=fmexbg,colframe=fmex,colbacktitle=fmexbg,coltitle=fmex,title={\faPencil*\ #1},
  titlerule=0.4pt,titlerule style={fmex!40,dashed}}
\newtcolorbox{fmexfill}[1]{fmbase,breakable=false,height fill,colback=fmexbg,colframe=fmex,colbacktitle=fmexbg,coltitle=fmex,title={\faPencil*\ #1},
  titlerule=0.4pt,titlerule style={fmex!40,dashed}}
\newtcolorbox{fmnote}{enhanced,breakable,colback=white,colframe=white,borderline west={3pt}{0pt}{fmaccent},arc=0pt,left=3mm,right=2mm,top=1mm,bottom=1mm,before skip=6pt,after skip=6pt,before upper=\raggedright}
\newcommand{\fmkey}[1]{\colorbox{fmyellow!45}{\strut #1}}
\newcommand{\fmuline}[1]{\textbf{\color{fmaccent}#1}}
\definecolor{fmok}{HTML}{15803D}\definecolor{fmev}{HTML}{2563EB}\definecolor{fmsf}{HTML}{C2410C}\definecolor{fmsi}{HTML}{7C3AED}\definecolor{fmas}{HTML}{A3203A}
\newcommand{\fmbadge}[2]{\tcbox[on line,colback=#1,colframe=#1,arc=2mm,boxsep=0pt,left=4pt,right=4pt,top=1.5pt,bottom=1.5pt]{\color{white}\footnotesize\bfseries #2}}
\newcommand{\fmpau}[1]{\tcbox[on line,colback=fmnavy,colframe=fmnavy,arc=2mm,boxsep=0pt,left=4pt,right=4pt,top=1.5pt,bottom=1.5pt]{\color{white}\scriptsize\bfseries\MakeUppercase{#1}}\ }
\newcommand{\fmchip}[1]{\tcbox[on line,colback=fmsf!10,colframe=fmsf!50,arc=2mm,boxsep=1pt,left=4pt,right=4pt]{#1}}
\newcommand{\fmlevel}[2]{\tcbox[on line,colback=white,colframe=fmmuted!40,arc=2mm,boxsep=2pt,halign=center]{\begin{tabular}{c}{\scriptsize\color{fmmuted}\MakeUppercase{#1}}\\#2\end{tabular}}}

% gràfiques
\pgfplotsset{fmaxis/.style={axis lines=middle,grid=major,grid style={fmmuted!15},axis line style={fmmuted,-latex},
  tick label style={font=\scriptsize,color=fmmuted},clip=true,unbounded coords=jump,restrict y to domain=-200:200,
  every axis plot/.append style={line join=round}},
  fmcurve/.style={fmaccent,line width=1.6pt,smooth}}
\tikzset{fmasy/.style={fmmuted,dashed,line width=0.9pt},fmguide/.style={fmaccent,dotted,line width=0.8pt},
  fmdot/.style={circle,fill=fmaccent,inner sep=0pt,minimum size=5pt},fmopen/.style={circle,draw=fmaccent,fill=white,line width=1pt,inner sep=0pt,minimum size=5pt}}

\begin{document}
\thispagestyle{primera}
\begin{tikzpicture}[remember picture,overlay]
  \fill[fmgreen] (current page.north west) rectangle ([yshift=-10.1cm]current page.north east);
  \begin{scope}
    \clip (current page.north west) rectangle ([yshift=-10.1cm]current page.north east);
    \node[anchor=north west,inner sep=0pt,opacity=0.15] at ([xshift=6.6cm,yshift=-1.75cm]current page.north west) {\fmcard};
  \end{scope}
  \node[anchor=north west,inner sep=0pt] at ([xshift=1.5cm,yshift=-1.1cm]current page.north west) {\tcbox[colback=white,colframe=white,arc=2.5mm,boxsep=0pt,left=7pt,right=7pt,top=5pt,bottom=5pt]{\includegraphics[height=0.75cm]{logo-fentmates.pdf}}};
\end{tikzpicture}
\vspace*{0.55cm}
\tcbox[colback=white!15!fmgreen,colframe=white!15!fmgreen,arc=2.5mm,boxsep=0pt,left=6pt,right=6pt,top=3pt,bottom=3pt]{\color{white}\small\bfseries Apuntes de teoría · 2.º Bachillerato CCSS II}

\vspace{4pt}
{\raggedright\color{white}\fontsize{25}{30}\selectfont\bfseries{\color{fmyellow}Unidad 10} · Distribución binomial y normal\par}
\vspace{6pt}
\begin{minipage}{\linewidth}\color{white!88!fmgreen}\small Teoría, ejemplos resueltos paso a paso: variables aleatorias discretas y continuas, la distribución binomial (con la fórmula y con la tabla de la PAU), la distribución normal, la normal tipificada \fmm{N(0,1)} y su tabla, la tipificación, los problemas inversos y la aproximación de la binomial por la normal. Incluye las dos tablas que se adjuntan al examen de la PAU valenciana, para que te acostumbres a usarlas.\end{minipage}
\vspace{0.7cm}

\begin{tcolorbox}[colback=fmsurface,colframe=fmsurface,arc=3mm,left=4mm,right=4mm,top=2mm,bottom=2mm]
\small\textbf{Índice}\par\smallskip
\begin{minipage}[t]{0.42\linewidth}\small \mbox{\hyperlink{sec1}{\textbf{\color{fmgreen}10.1}\ \ Variables aleatorias discretas}}\\[3pt]
\mbox{\hyperlink{sec2}{\textbf{\color{fmgreen}10.2}\ \ La distribución binomial}}\\[3pt]
\mbox{\hyperlink{sec3}{\textbf{\color{fmgreen}10.3}\ \ La tabla de la binomial}}\\[3pt]
\mbox{\hyperlink{sec4}{\textbf{\color{fmgreen}10.4}\ \ Variables continuas}}\\[3pt]\end{minipage}\hfill\begin{minipage}[t]{0.56\linewidth}\small \mbox{\hyperlink{sec5}{\textbf{\color{fmgreen}10.5}\ \ La distribución normal}}\\[3pt]
\mbox{\hyperlink{sec6}{\textbf{\color{fmgreen}10.6}\ \ La normal tipificada y su tabla}}\\[3pt]
\mbox{\hyperlink{sec7}{\textbf{\color{fmgreen}10.7}\ \ Tipificación. Problemas}}\\[3pt]
\mbox{\hyperlink{sec8}{\textbf{\color{fmgreen}10.8}\ \ Aproximación de la binomial por la normal}}\\[3pt]
\mbox{\hyperlink{annex}{\textbf{\color{fmgreen}Anexo}\ \ Tablas de la PAU}}\\[3pt]\end{minipage}
\end{tcolorbox}

\Needspace{18\baselineskip}
\section{Variables aleatorias discretas}\hypertarget{sec1}{}

  

  
\Needspace{7\baselineskip}
\begin{fmdef}{Variable aleatoria}
Una \textbf{variable aleatoria} \fmm{X} asigna un número a cada resultado de un experimento aleatorio. Por ejemplo, al lanzar tres monedas, \fmm{X=}«número de caras» toma los valores 0, 1, 2 y 3.

      
\begin{itemize}
\item Es \textbf{discreta} si toma valores aislados (que se pueden contar): número de caras, de clientes, de piezas defectuosas\dots{}
\item Es \textbf{continua} si puede tomar cualquier valor de un intervalo: altura, peso, tiempo, nota\dots{}
\end{itemize}
\end{fmdef}

  
\Needspace{7\baselineskip}
\begin{fmdef}{Función de probabilidad, media y desviación típica}
La \textbf{función de probabilidad} de una variable discreta da \fmm{p_i=P(X=x_i)} para cada valor. Se cumple \fmm{p_i\geq0} y \fmm{\sum p_i=1}. La \textbf{media} (o esperanza) y la \textbf{varianza} son:

      
\[\mu=E(X)=\sum x_i\,p_i\qquad\qquad\sigma^2=\sum x_i^2\,p_i-\mu^2\qquad\qquad\sigma=\sqrt{\sigma^2}\]

      

La media es el valor que esperamos obtener «de media» si repetimos el experimento muchas veces; la desviación típica mide cómo de dispersos están los valores.
\end{fmdef}

  
\Needspace{7\baselineskip}
\begin{fmex}{Ejemplo resuelto}
\begin{enumerate}[label=\textbf{\arabic*.},leftmargin=*]
\item Lanzamos tres monedas y \fmm{X=}«número de caras». Halla la función de probabilidad, la media y la desviación típica.
          

            

De los 8 resultados equiprobables, hay 1, 3, 3 y 1 con 0, 1, 2 y 3 caras:

              
\begin{center}\renewcommand{\arraystretch}{1.5}\begin{tabular}{ccccc}\toprule
$x_i$ & 0 & 1 & 2 & 3\\
$p_i$ & $\frac18$ & $\frac38$ & $\frac38$ & $\frac18$\\
\bottomrule\end{tabular}\end{center}

            

\fmm{\mu=0\cdot\frac18+1\cdot\frac38+2\cdot\frac38+3\cdot\frac18=\frac{12}{8}=1{,}5} caras

            

\fmm{\sigma^2=0+\frac38+4\cdot\frac38+9\cdot\frac18-1{,}5^2=3-2{,}25=0{,}75\Rightarrow\sigma=\sqrt{0{,}75}\approx0{,}866}
\end{enumerate}
\end{fmex}

\Needspace{18\baselineskip}
\section{La distribución binomial}\hypertarget{sec2}{}

  

  
\Needspace{7\baselineskip}
\begin{fmdef}{¿Cuándo es binomial?}
Una variable \fmm{X} sigue una \textbf{distribución binomial} \fmm{B(n,p)} cuando cuenta el \textbf{número de éxitos} en un experimento que cumple:

      
\begin{enumerate}
\item Se repite una misma prueba \fmm{n} veces.
\item Cada prueba solo tiene dos resultados: \textbf{éxito} (\fmm{A}) o \textbf{fracaso} (\fmm{\overline A}).
\item La probabilidad de éxito, \fmm{p=P(A)}, es la misma en todas las pruebas (la de fracaso es \fmm{q=1-p}).
\item Las pruebas son \textbf{independientes} (por ejemplo, extracciones \textbf{con} reemplazamiento, o de una población muy grande).
\end{enumerate}

      
\begin{fmnote}Ejemplos: número de caras en 10 lanzamientos, número de piezas defectuosas en una muestra de 20 (si el 3 \% lo son), número de aciertos en un test contestado al azar\dots{} Cuidado: las extracciones \textbf{sin} reemplazamiento de una urna pequeña \textbf{no} son binomiales (la \fmm{p} cambia).\end{fmnote}
\end{fmdef}

  
\Needspace{7\baselineskip}
\begin{fmprop}{Fórmula de la binomial}

\[P(X=k)=\binom{n}{k}\,p^k\,q^{\,n-k},\qquad k=0,1,\dots,n\qquad\qquad\binom nk=\frac{n!}{k!\,(n-k)!}\]

      

El \textbf{número combinatorio} \fmm{\binom nk} («\fmm{n} sobre \fmm{k}») cuenta de cuántas maneras se pueden colocar los \fmm{k} éxitos entre las \fmm{n} pruebas; \fmm{p^kq^{n-k}} es la probabilidad de cada una de esas maneras. \\
Recuerda: \fmm{\binom n0=\binom nn=1}, \fmm{\binom n1=n}, \fmm{\binom nk=\binom n{n-k}}. En la calculadora, tecla \textbf{nCr}.

      
\[\mu=n\,p\qquad\qquad\sigma=\sqrt{n\,p\,q}\]

      

Para probabilidades acumuladas sumamos: \fmm{P(X\leq k)=P(X=0)+\dots+P(X=k)}. Y cuando «al menos uno» o «más de\dots{}» tenga muchos términos, usamos el contrario: \fmm{P(X\geq1)=1-P(X=0)}.
\end{fmprop}

  
\Needspace{7\baselineskip}
\begin{fmex}{Ejemplos resueltos}
\begin{enumerate}[label=\textbf{\arabic*.},leftmargin=*]
\item Un test tiene 10 preguntas con 4 opciones cada una. Un alumno contesta todas al azar. Probabilidad de que acierte exactamente 3, y de que acierte al menos 1.
          

            

\fmm{X=}«número de aciertos» es \fmm{B(10;\,0{,}25)}: 10 pruebas independientes con \fmm{p=\frac14} de acierto.

            

\fmm{P(X=3)=\binom{10}{3}\,0{,}25^3\,0{,}75^7=120\cdot0{,}015625\cdot0{,}1335\approx0{,}2503}

            

\fmm{P(X\geq1)=1-P(X=0)=1-0{,}75^{10}\approx1-0{,}0563=0{,}9437}

            

Media: \fmm{\mu=10\cdot0{,}25=2{,}5} aciertos; \fmm{\sigma=\sqrt{10\cdot0{,}25\cdot0{,}75}\approx1{,}37}.
\item Un 48 \% de los habitantes de una ciudad son hombres. Elegimos 4 personas con reemplazamiento y \fmm{X} cuenta los hombres. ¿Qué distribución sigue \fmm{X}? Calcula \fmm{P(X=2)}. \fmpau{PAU Aragón 2024}
          

            

\fmm{X\sim B(4;\,0{,}48)}; \fmm{\ P(X=2)=\binom42\,0{,}48^2\,0{,}52^2=6\cdot0{,}2304\cdot0{,}2704\approx0{,}3738}
\end{enumerate}
\end{fmex}

  
\Needspace{18\baselineskip}
\section{La tabla de la binomial de la PAU}\hypertarget{sec3}{}

  

  
\Needspace{7\baselineskip}
\begin{fmdef}{Cómo se lee la tabla}
En la PAU valenciana se adjunta esta tabla. Da directamente \fmm{P(X=x)} para \fmm{n=2,\dots,9} y los valores de \fmm{p} de la cabecera. Buscamos el bloque de la \fmm{n}, la fila de la \fmm{x} y la columna de la \fmm{p}.

      
\begin{itemize}
\item \textbf{Probabilidades acumuladas}: sumamos varias casillas de la misma columna, o hacemos el contrario.
\item \textbf{Si \fmm{p\gt0{,}5}} (no está en la tabla): contamos los fracasos. Si \fmm{X\sim B(n,p)}, entonces \fmm{Y=n-X\sim B(n,q)} con \fmm{q=1-p\lt0{,}5}, y \fmm{P(X=k)=P(Y=n-k)}.
\item \textbf{Si \fmm{n} o \fmm{p} no están en la tabla} (\fmm{n\geq10}, \fmm{p=0{,}62}\dots{}): usamos la fórmula; y si \fmm{n} es grande, la aproximación por la normal (apartado 10.8).
\end{itemize}
\end{fmdef}

  
\begin{fmnote}\textbf{La tabla de la PAU.} La tabla completa, la misma que se adjunta al examen de selectividad, se encuentra en el anexo (página \pageref{tbin}).\end{fmnote}

  
\Needspace{7\baselineskip}
\begin{fmex}{Ejemplos resueltos con la tabla}
\begin{enumerate}[label=\textbf{\arabic*.},leftmargin=*]
\item Un 20 \% de los clientes de una tienda paga con el móvil. Elegimos 8 clientes. Probabilidad de que como mucho 2 paguen con el móvil.
          

            

\fmm{X\sim B(8;\,0{,}20)}. Columna \fmm{p=0{,}20}, bloque \fmm{n=8}:

            

\fmm{P(X\leq2)=P(0)+P(1)+P(2)=0{,}1678+0{,}3355+0{,}2936=0{,}7969}
\item El 70 \% de los alumnos aprueba. Elegimos 6 alumnos. Probabilidad de que aprueben exactamente 5, y de que aprueben al menos 5.
          

            

\fmm{p=0{,}7} no está en la tabla. Contamos los suspensos: \fmm{Y\sim B(6;\,0{,}30)}, y «aprueban 5» es «suspende 1».

            

\fmm{P(X=5)=P(Y=1)=0{,}3025}

            

\fmm{P(X\geq5)=P(Y\leq1)=P(Y=0)+P(Y=1)=0{,}1176+0{,}3025=0{,}4201}
\item El 87 \% de los jóvenes compra por Internet. En una encuesta a 10 jóvenes, probabilidad de que exactamente 2 no compren. \fmpau{PAU Castilla y León 2026 (adaptado)}
          

            

\fmm{X=}«no compran» \fmm{\sim B(10;\,0{,}13)}: ni \fmm{n=10} ni \fmm{p=0{,}13} están en la tabla. Fórmula:

            

\fmm{P(X=2)=\binom{10}2\,0{,}13^2\,0{,}87^8=45\cdot0{,}0169\cdot0{,}3282\approx0{,}2496}
\end{enumerate}
\end{fmex}

\Needspace{18\baselineskip}
\section{Variables aleatorias continuas}\hypertarget{sec4}{}

  

  
\Needspace{7\baselineskip}
\begin{fmdef}{Función de densidad}
\par\noindent \begin{minipage}[t]{0.48\linewidth}\vspace{0pt}\raggedright Una variable continua puede tomar infinitos valores, y la probabilidad de cada valor concreto es 0. En lugar de barras, la describe una \textbf{función de densidad} \fmm{f(x)\geq0}, con área total bajo la curva igual a 1.

          
\[\begin{aligned}P(a\leq X\leq b)&=\int_a^b f(x)\,dx\\&=\text{área bajo }f\text{ entre }a\text{ y }b\end{aligned}\]

          
\begin{fmnote}Como \fmm{P(X=a)=0}, en una variable continua \textbf{da igual} poner \fmm{\lt} o \fmm{\leq}: \fmm{P(X\lt a)=P(X\leq a)}. En la binomial (discreta), en cambio, ¡\textbf{no} es lo mismo!\end{fmnote}\end{minipage}\hfill\begin{minipage}[t]{0.48\linewidth}\vspace{0pt}\raggedright {\centering \begin{tikzpicture}[x=0.8235cm,y=2.4000cm]\fill[fmaccent,fill opacity=0.35] (-0.4,0) -- plot[domain=-0.4:1.3,samples=70,smooth] (\x,{exp(-\x*\x/2)}) -- (1.3,0) -- cycle;\draw[fmmuted,line width=0.7pt] (-3.4,0) -- (3.4,0);\draw[fmgreen,line width=1.4pt] plot[domain=-3.4:3.4,samples=120,smooth] (\x,{exp(-\x*\x/2)});\draw[fmmuted,line width=0.7pt] (-0.4,0) -- (-0.4,-0.06) node[below,font=\footnotesize,inner sep=2pt,text=black] {$a$};\draw[fmmuted,line width=0.7pt] (1.3,0) -- (1.3,-0.06) node[below,font=\footnotesize,inner sep=2pt,text=black] {$b$};\node[anchor=north east,font=\footnotesize] at (3.4,1.08) {$y=f(x)$};\end{tikzpicture}\par}{\centering\footnotesize La probabilidad es el área sombreada\par}\end{minipage}\par\medskip
\end{fmdef}

\Needspace{18\baselineskip}
\section{La distribución normal $N(\mu,\sigma)$}\hypertarget{sec5}{}

  

  
\Needspace{7\baselineskip}
\begin{fmdef}{La campana de Gauss}
Muchas variables reales (alturas, pesos, notas, tiempos, errores de medida\dots{}) siguen una \textbf{distribución normal} de media \fmm{\mu} y desviación típica \fmm{\sigma}, \fmm{X\sim N(\mu,\sigma)}. Su densidad es la \textbf{campana de Gauss}:

      
\[f(x)=\frac{1}{\sigma\sqrt{2\pi}}\,e^{-\frac{(x-\mu)^2}{2\sigma^2}}\]

      
\begin{itemize}
\item Es \textbf{simétrica} respecto de \fmm{x=\mu}, donde tiene el máximo: \fmm{P(X\leq\mu)=P(X\geq\mu)=0{,}5}.
\item \fmm{\mu} indica dónde está centrada; \fmm{\sigma}, cómo de ancha es: cuanto mayor \fmm{\sigma}, más baja y ancha.
\item Regla 68--95--99,7: entre \fmm{\mu\pm\sigma} hay aproximadamente el 68 \% de los valores; entre \fmm{\mu\pm2\sigma}, el 95 \%; entre \fmm{\mu\pm3\sigma}, el 99,7 \%.
\end{itemize}

      
\begin{fmnote}Nunca integraremos la función de densidad (no tiene primitiva elemental): las probabilidades se calculan con la \textbf{tabla de la normal tipificada}.\end{fmnote}
\end{fmdef}

\Needspace{18\baselineskip}
\section{La normal tipificada $N(0,1)$ y su tabla}\hypertarget{sec6}{}

  

  
\Needspace{7\baselineskip}
\begin{fmdef}{Cómo se lee la tabla de la PAU}
La \textbf{normal tipificada} es \fmm{Z\sim N(0,1)}. La tabla de la PAU da \fmm{F(z)=P(Z\leq z)}, el área a la \textbf{izquierda} de \fmm{z}, para \fmm{z} entre 0,00 y 3,99. La fila da las unidades y las décimas de \fmm{z} y la columna, las centésimas. Por ejemplo, \fmm{P(Z\leq1{,}25)}: fila 1,2 y columna 0,05 \fmm{\to0{,}8944}.

      

Con la simetría y el contrario se calcula todo:

      
\begin{center}\begin{minipage}[t]{0.32\linewidth}\centering {\centering \begin{tikzpicture}[x=0.5294cm,y=1.5500cm]\fill[fmaccent,fill opacity=0.35] (-3.4,0) -- plot[domain=-3.4:0.8,samples=70,smooth] (\x,{exp(-\x*\x/2)}) -- (0.8,0) -- cycle;\draw[fmmuted,line width=0.7pt] (-3.4,0) -- (3.4,0);\draw[fmgreen,line width=1.4pt] plot[domain=-3.4:3.4,samples=120,smooth] (\x,{exp(-\x*\x/2)});\draw[fmmuted,line width=0.7pt] (0,0) -- (0,-0.06) node[below,font=\footnotesize,inner sep=2pt,text=black] {$0$};\draw[fmmuted,line width=0.7pt] (0.8,0) -- (0.8,-0.06) node[below,font=\footnotesize,inner sep=2pt,text=black] {$z$};\end{tikzpicture}\par}{\centering\footnotesize \fmm{P(Z\leq z)}: tabla\par}\end{minipage}\hfill\begin{minipage}[t]{0.32\linewidth}\centering {\centering \begin{tikzpicture}[x=0.5294cm,y=1.5500cm]\fill[fmaccent,fill opacity=0.35] (0.8,0) -- plot[domain=0.8:3.4,samples=70,smooth] (\x,{exp(-\x*\x/2)}) -- (3.4,0) -- cycle;\draw[fmmuted,line width=0.7pt] (-3.4,0) -- (3.4,0);\draw[fmgreen,line width=1.4pt] plot[domain=-3.4:3.4,samples=120,smooth] (\x,{exp(-\x*\x/2)});\draw[fmmuted,line width=0.7pt] (0,0) -- (0,-0.06) node[below,font=\footnotesize,inner sep=2pt,text=black] {$0$};\draw[fmmuted,line width=0.7pt] (0.8,0) -- (0.8,-0.06) node[below,font=\footnotesize,inner sep=2pt,text=black] {$z$};\end{tikzpicture}\par}{\centering\footnotesize \fmm{P(Z\geq z)=1-P(Z\leq z)}\par}\end{minipage}\hfill\begin{minipage}[t]{0.32\linewidth}\centering {\centering \begin{tikzpicture}[x=0.5294cm,y=1.5500cm]\fill[fmaccent,fill opacity=0.35] (-3.4,0) -- plot[domain=-3.4:-0.8,samples=70,smooth] (\x,{exp(-\x*\x/2)}) -- (-0.8,0) -- cycle;\draw[fmmuted,line width=0.7pt] (-3.4,0) -- (3.4,0);\draw[fmgreen,line width=1.4pt] plot[domain=-3.4:3.4,samples=120,smooth] (\x,{exp(-\x*\x/2)});\draw[fmmuted,line width=0.7pt] (-0.8,0) -- (-0.8,-0.06) node[below,font=\footnotesize,inner sep=2pt,text=black] {$-z$};\draw[fmmuted,line width=0.7pt] (0,0) -- (0,-0.06) node[below,font=\footnotesize,inner sep=2pt,text=black] {$0$};\end{tikzpicture}\par}{\centering\footnotesize \fmm{P(Z\leq -z)=1-P(Z\leq z)}\par}\end{minipage}\\[10pt]\begin{minipage}[t]{0.32\linewidth}\centering {\centering \begin{tikzpicture}[x=0.5294cm,y=1.5500cm]\fill[fmaccent,fill opacity=0.35] (-0.8,0) -- plot[domain=-0.8:3.4,samples=70,smooth] (\x,{exp(-\x*\x/2)}) -- (3.4,0) -- cycle;\draw[fmmuted,line width=0.7pt] (-3.4,0) -- (3.4,0);\draw[fmgreen,line width=1.4pt] plot[domain=-3.4:3.4,samples=120,smooth] (\x,{exp(-\x*\x/2)});\draw[fmmuted,line width=0.7pt] (-0.8,0) -- (-0.8,-0.06) node[below,font=\footnotesize,inner sep=2pt,text=black] {$-z$};\draw[fmmuted,line width=0.7pt] (0,0) -- (0,-0.06) node[below,font=\footnotesize,inner sep=2pt,text=black] {$0$};\end{tikzpicture}\par}{\centering\footnotesize \fmm{P(Z\geq -z)=P(Z\leq z)}\par}\end{minipage}\hfill\begin{minipage}[t]{0.32\linewidth}\centering {\centering \begin{tikzpicture}[x=0.5294cm,y=1.5500cm]\fill[fmaccent,fill opacity=0.35] (-0.6,0) -- plot[domain=-0.6:1.4,samples=70,smooth] (\x,{exp(-\x*\x/2)}) -- (1.4,0) -- cycle;\draw[fmmuted,line width=0.7pt] (-3.4,0) -- (3.4,0);\draw[fmgreen,line width=1.4pt] plot[domain=-3.4:3.4,samples=120,smooth] (\x,{exp(-\x*\x/2)});\draw[fmmuted,line width=0.7pt] (-0.6,0) -- (-0.6,-0.06) node[below,font=\footnotesize,inner sep=2pt,text=black] {$a$};\draw[fmmuted,line width=0.7pt] (1.4,0) -- (1.4,-0.06) node[below,font=\footnotesize,inner sep=2pt,text=black] {$b$};\end{tikzpicture}\par}{\centering\footnotesize \fmm{P(a\leq Z\leq b)=P(Z\leq b)-P(Z\leq a)}\par}\end{minipage}\hfill\begin{minipage}[t]{0.32\linewidth}\centering {\centering \begin{tikzpicture}[x=0.5294cm,y=1.5500cm]\fill[fmaccent,fill opacity=0.35] (-1.2,0) -- plot[domain=-1.2:1.2,samples=70,smooth] (\x,{exp(-\x*\x/2)}) -- (1.2,0) -- cycle;\draw[fmmuted,line width=0.7pt] (-3.4,0) -- (3.4,0);\draw[fmgreen,line width=1.4pt] plot[domain=-3.4:3.4,samples=120,smooth] (\x,{exp(-\x*\x/2)});\draw[fmmuted,line width=0.7pt] (-1.2,0) -- (-1.2,-0.06) node[below,font=\footnotesize,inner sep=2pt,text=black] {$-z$};\draw[fmmuted,line width=0.7pt] (1.2,0) -- (1.2,-0.06) node[below,font=\footnotesize,inner sep=2pt,text=black] {$z$};\end{tikzpicture}\par}{\centering\footnotesize \fmm{P(-z\leq Z\leq z)=2P(Z\leq z)-1}\par}\end{minipage}\end{center}

      
\begin{fmnote}\fmuline{Problema inverso}: si nos dan la probabilidad y queremos \fmm{z}, la buscamos \textbf{dentro} de la tabla y leemos la fila y la columna (si no está exacta, tomamos la más próxima). Si la probabilidad es menor que 0,5, \fmm{z} es negativo: buscamos \fmm{1-p} y cambiamos el signo.\end{fmnote}
\end{fmdef}

  
\begin{fmnote}\textbf{La tabla de la PAU.} La tabla completa, la misma que se adjunta al examen de selectividad, se encuentra en el anexo (página \pageref{tnorm}).\end{fmnote}

  
\Needspace{7\baselineskip}
\begin{fmex}{Ejemplos resueltos}
\begin{enumerate}[label=\textbf{\arabic*.},leftmargin=*]
\item Si \fmm{Z\sim N(0,1)}, calcula \fmm{P(Z\gt1{,}25)}, \fmm{P(Z\lt-0{,}84)} y \fmm{P(-1\lt Z\lt2)}.
          

            

\fmm{P(Z\gt1{,}25)=1-P(Z\leq1{,}25)=1-0{,}8944=0{,}1056}

            

Por simetría: \fmm{P(Z\lt-0{,}84)=P(Z\gt0{,}84)=1-0{,}7995=0{,}2005}

            

\fmm{P(-1\lt Z\lt2)=P(Z\lt2)-P(Z\lt-1)=0{,}9772-(1-0{,}8413)=0{,}8185}
\item Halla \fmm{k} tal que \fmm{P(Z\leq k)=0{,}9} y \fmm{c} tal que \fmm{P(Z\geq c)=0{,}95}.
          

            

Dentro de la tabla, el valor más próximo a 0,9 es 0,8997 (fila 1,2, columna 0,08): \fmm{k\approx1{,}28}.

            

\fmm{P(Z\geq c)=0{,}95\gt0{,}5}: \fmm{c} es negativo y, por simetría, \fmm{P(Z\leq-c)=0{,}95}. En la tabla, 0,9495 (1,64) y 0,9505 (1,65) están igual de cerca: \fmm{-c\approx1{,}645\Rightarrow c\approx-1{,}645}.
\end{enumerate}
\end{fmex}

  
\Needspace{18\baselineskip}
\section{Tipificación. Problemas con la normal}\hypertarget{sec7}{}

  

  
\Needspace{7\baselineskip}
\begin{fmprop}{Tipificación}
Si \fmm{X\sim N(\mu,\sigma)}, la variable \fmm{Z=\dfrac{X-\mu}{\sigma}} sigue una \fmm{N(0,1)}. Por tanto:

      
\[P(X\leq a)=P\!\left(Z\leq\frac{a-\mu}{\sigma}\right)\qquad\qquad P(a\leq X\leq b)=P\!\left(\frac{a-\mu}{\sigma}\leq Z\leq\frac{b-\mu}{\sigma}\right)\]

      

El valor \fmm{z=\frac{a-\mu}{\sigma}} dice a cuántas desviaciones típicas está \fmm{a} de la media. En la PAU redondeamos \fmm{z} a dos decimales para poder usar la tabla.

      
\begin{fmnote}\fmuline{Tipos de problemas}: (1) \textbf{directos}: nos dan \fmm{a} y queremos la probabilidad; (2) \textbf{«cuántas personas\dots{}»}: multiplicamos la probabilidad por el total; (3) \textbf{inversos}: nos dan la probabilidad y queremos el valor \fmm{a} (\fmm{a=\mu+z\sigma}), o bien \fmm{\mu} o \fmm{\sigma}.\end{fmnote}
\end{fmprop}

  
\Needspace{7\baselineskip}
\begin{fmex}{Ejemplos resueltos}
\begin{enumerate}[label=\textbf{\arabic*.},leftmargin=*]
\item El coeficiente intelectual de una población de 1200 personas sigue una \fmm{N(100,14)}. ¿Cuántas personas tienen un CI entre 108 y 116? ¿Y mayor de 120? \fmpau{PAU C. Valenciana 2001}
          

            

\[\begin{aligned}P(108\leq X\leq116)&=P\!\left(\frac{108-100}{14}\leq Z\leq\frac{116-100}{14}\right)=P(0{,}57\leq Z\leq1{,}14)\\&=0{,}8729-0{,}7157=0{,}1572\end{aligned}\]

            

\fmm{1200\cdot0{,}1572\approx188{,}6\to} unas 189 personas.

            

\fmm{P(X\gt120)=P(Z\gt1{,}43)=1-0{,}9236=0{,}0764\Rightarrow1200\cdot0{,}0764\approx92} personas.
\item El tiempo que tardan los alumnos en hacer un examen sigue una \fmm{N(60,10)} (minutos). ¿Cuánto tiempo hay que dar para que lo acabe el 96 \% de los alumnos? \fmpau{PAU País Vasco 2020}
          

            

Buscamos \fmm{t} con \fmm{P(X\leq t)=0{,}96}. En la tabla, el más próximo a 0,96 es 0,9599 (\fmm{z=1{,}75}).

            

\fmm{\frac{t-60}{10}=1{,}75\Rightarrow t=60+1{,}75\cdot10=77{,}5} minutos.
\item La nota de un examen sigue una normal de media 5,8. El 35 \% de los alumnos saca más de 6,8. Calcula la desviación típica. \fmpau{PAU País Vasco 2024}
          

            

\fmm{P(X\gt6{,}8)=0{,}35\Rightarrow P\!\left(Z\leq\frac{1}{\sigma}\right)=0{,}65}. En la tabla, 0,6517 (\fmm{z=0{,}39}) es el más próximo.

            

\fmm{\frac{1}{\sigma}=0{,}39\Rightarrow\sigma=\frac{1}{0{,}39}\approx2{,}56} puntos.
\end{enumerate}
\end{fmex}

  
\Needspace{18\baselineskip}
\section{Aproximación de la binomial por la normal}\hypertarget{sec8}{}

  

  
\Needspace{7\baselineskip}
\begin{fmprop}{Teorema de De Moivre-Laplace}
Cuando \fmm{n} es grande, calcular la binomial con la fórmula es muy pesado (¡\fmm{P(X\geq40)} con \fmm{n=100} tendría 61 términos!). Si \fmm{np\geq5} y \fmm{nq\geq5}, la binomial se parece mucho a una normal con la misma media y desviación típica:

      
\[X\sim B(n,p)\ \approx\ X'\sim N\!\left(np,\ \sqrt{npq}\right)\]

      

\textbf{Corrección de continuidad}: cada valor entero \fmm{k} de la binomial es una barra que va de \fmm{k-0{,}5} a \fmm{k+0{,}5}. Por eso:

      
\begin{center}\renewcommand{\arraystretch}{1.7}\small\begin{tabularx}{\linewidth}{>{\raggedright\arraybackslash}p{5.2cm}>{\centering\arraybackslash}X}\toprule
\textbf{Binomial} & \textbf{Normal}\\\midrule
$P(X=k)$ & $P(k-0{,}5\leq X'\leq k+0{,}5)$\\[3pt]
$P(X\leq k)$ & $P(X'\leq k+0{,}5)$\\[3pt]
$P(X\lt k)=P(X\leq k-1)$ & $P(X'\leq k-0{,}5)$\\[3pt]
$P(X\geq k)$ & $P(X'\geq k-0{,}5)$\\[3pt]
$P(X\gt k)=P(X\geq k+1)$ & $P(X'\geq k+0{,}5)$\\[3pt]
$P(a\leq X\leq b)$ & $P(a-0{,}5\leq X'\leq b+0{,}5)$\\[3pt]
\bottomrule\end{tabularx}\end{center}

      
\begin{fmnote}Truco para no equivocarse: escribe primero los enteros que \textbf{sí} entran (por ejemplo, «más de 499 y como mucho 580» son 500, \dots{}, 580) y después alarga medio punto hacia cada lado: \fmm{499{,}5\leq X'\leq580{,}5}.\end{fmnote}
\end{fmprop}

  
\Needspace{7\baselineskip}
\begin{fmex}{Ejemplos resueltos}
\begin{enumerate}[label=\textbf{\arabic*.},leftmargin=*]
\item Un creador de vídeos con talento consigue que un vídeo se haga viral con probabilidad 0,5. Probabilidad de que, de los próximos 100 vídeos, al menos 40 se hagan virales. \fmpau{PAU C. Valenciana 2026}
          

            

\fmm{X\sim B(100;\,0{,}5)}: \fmm{np=nq=50\geq5}. \fmm{\mu=50}, \fmm{\sigma=\sqrt{100\cdot0{,}5\cdot0{,}5}=5\Rightarrow X'\sim N(50,5)}.

            

\fmm{P(X\geq40)\approx P(X'\geq39{,}5)=P\!\left(Z\geq\frac{39{,}5-50}{5}\right)=P(Z\geq-2{,}1)=P(Z\leq2{,}1)=0{,}9821}
\item Una aseguradora tiene 2500 coches asegurados y cada uno tiene un accidente en un año con probabilidad 0,1. Probabilidad de que al menos 270 tengan accidente. \fmpau{PAU Canarias 2022}
          

            

\fmm{X\sim B(2500;\,0{,}1)\approx N(250,\ \sqrt{225}=15)}

            

\fmm{P(X\geq270)\approx P(X'\geq269{,}5)=P(Z\geq1{,}3)=1-0{,}9032=0{,}0968}
\end{enumerate}
\end{fmex}

  
\clearpage
\hypertarget{annex}{}\phantomsection
{\color{fmnavy}\Large\bfseries\colorbox{fmdef}{\color{white}\strut\ Anexo\ }\hspace{0.6em}Tablas de la PAU\par}\vspace{2pt}{\color{fmdef}\titlerule[1.2pt]}
\addcontentsline{toc}{section}{Anexo. Tablas de la PAU}
\vspace{6pt}
\label{tbin}{\color{fmgreen}\large\bfseries Tabla de la distribución binomial}\par\smallskip
Da \fmm{P[X=x]=\binom nx\,p^x(1-p)^{n-x}} para \fmm{n=2,\dots,9}.\newline Si \fmm{p\gt0{,}5}, cuenta los fracasos: \fmm{Y=n-X\sim B(n,1-p)} (apartado 10.3).\par\smallskip
\begin{center}\linespread{1}\selectfont\scriptsize\setlength{\tabcolsep}{4.2pt}\renewcommand{\arraystretch}{1.0}
\begin{tabular}{cc|ccccccccccccc}\toprule
\multirow{2}{*}{$n$} & \multirow{2}{*}{$x$} & \multicolumn{13}{c}{$p$}\\
 & & \textbf{0,01} & \textbf{0,05} & \textbf{0,10} & \textbf{0,15} & \textbf{0,20} & \textbf{0,25} & \textbf{0,30} & \textbf{1/3} & \textbf{0,35} & \textbf{0,40} & \textbf{0,45} & \textbf{0,49} & \textbf{0,50}\\
\midrule
\multirow{3}{*}{\textbf{2}} & \textbf{0} & 0,9801 & 0,9025 & 0,8100 & 0,7225 & 0,6400 & 0,5625 & 0,4900 & 0,4444 & 0,4225 & 0,3600 & 0,3025 & 0,2601 & 0,2500\\
 & \textbf{1} & 0,0198 & 0,0950 & 0,1800 & 0,2550 & 0,3200 & 0,3750 & 0,4200 & 0,4444 & 0,4550 & 0,4800 & 0,4950 & 0,4998 & 0,5000\\
 & \textbf{2} & 0,0001 & 0,0025 & 0,0100 & 0,0225 & 0,0400 & 0,0625 & 0,0900 & 0,1111 & 0,1225 & 0,1600 & 0,2025 & 0,2401 & 0,2500\\
\midrule
\multirow{4}{*}{\textbf{3}} & \textbf{0} & 0,9703 & 0,8574 & 0,7290 & 0,6141 & 0,5120 & 0,4219 & 0,3430 & 0,2963 & 0,2746 & 0,2160 & 0,1664 & 0,1327 & 0,1250\\
 & \textbf{1} & 0,0294 & 0,1354 & 0,2430 & 0,3251 & 0,3840 & 0,4219 & 0,4410 & 0,4444 & 0,4436 & 0,4320 & 0,4084 & 0,3823 & 0,3750\\
 & \textbf{2} & 0,0003 & 0,0071 & 0,0270 & 0,0574 & 0,0960 & 0,1406 & 0,1890 & 0,2222 & 0,2389 & 0,2880 & 0,3341 & 0,3674 & 0,3750\\
 & \textbf{3} & 0,0000 & 0,0001 & 0,0010 & 0,0034 & 0,0080 & 0,0156 & 0,0270 & 0,0370 & 0,0429 & 0,0640 & 0,0911 & 0,1176 & 0,1250\\
\midrule
\multirow{5}{*}{\textbf{4}} & \textbf{0} & 0,9606 & 0,8145 & 0,6561 & 0,5220 & 0,4096 & 0,3164 & 0,2401 & 0,1975 & 0,1785 & 0,1296 & 0,0915 & 0,0677 & 0,0625\\
 & \textbf{1} & 0,0388 & 0,1715 & 0,2916 & 0,3685 & 0,4096 & 0,4219 & 0,4116 & 0,3951 & 0,3845 & 0,3456 & 0,2995 & 0,2600 & 0,2500\\
 & \textbf{2} & 0,0006 & 0,0135 & 0,0486 & 0,0975 & 0,1536 & 0,2109 & 0,2646 & 0,2963 & 0,3105 & 0,3456 & 0,3675 & 0,3747 & 0,3750\\
 & \textbf{3} & 0,0000 & 0,0005 & 0,0036 & 0,0115 & 0,0256 & 0,0469 & 0,0756 & 0,0988 & 0,1115 & 0,1536 & 0,2005 & 0,2400 & 0,2500\\
 & \textbf{4} & 0,0000 & 0,0000 & 0,0001 & 0,0005 & 0,0016 & 0,0039 & 0,0081 & 0,0123 & 0,0150 & 0,0256 & 0,0410 & 0,0576 & 0,0625\\
\midrule
\multirow{6}{*}{\textbf{5}} & \textbf{0} & 0,9510 & 0,7738 & 0,5905 & 0,4437 & 0,3277 & 0,2373 & 0,1681 & 0,1317 & 0,1160 & 0,0778 & 0,0503 & 0,0345 & 0,0312\\
 & \textbf{1} & 0,0480 & 0,2036 & 0,3281 & 0,3915 & 0,4096 & 0,3955 & 0,3601 & 0,3292 & 0,3124 & 0,2592 & 0,2059 & 0,1657 & 0,1562\\
 & \textbf{2} & 0,0010 & 0,0214 & 0,0729 & 0,1382 & 0,2048 & 0,2637 & 0,3087 & 0,3292 & 0,3364 & 0,3456 & 0,3369 & 0,3185 & 0,3125\\
 & \textbf{3} & 0,0000 & 0,0011 & 0,0081 & 0,0244 & 0,0512 & 0,0879 & 0,1323 & 0,1646 & 0,1811 & 0,2304 & 0,2757 & 0,3060 & 0,3125\\
 & \textbf{4} & 0,0000 & 0,0000 & 0,0005 & 0,0022 & 0,0064 & 0,0146 & 0,0283 & 0,0412 & 0,0488 & 0,0768 & 0,1128 & 0,1470 & 0,1562\\
 & \textbf{5} & 0,0000 & 0,0000 & 0,0000 & 0,0001 & 0,0003 & 0,0010 & 0,0024 & 0,0041 & 0,0053 & 0,0102 & 0,0185 & 0,0282 & 0,0312\\
\midrule
\multirow{7}{*}{\textbf{6}} & \textbf{0} & 0,9415 & 0,7351 & 0,5314 & 0,3771 & 0,2621 & 0,1780 & 0,1176 & 0,0878 & 0,0754 & 0,0467 & 0,0277 & 0,0176 & 0,0156\\
 & \textbf{1} & 0,0571 & 0,2321 & 0,3543 & 0,3993 & 0,3932 & 0,3560 & 0,3025 & 0,2634 & 0,2437 & 0,1866 & 0,1359 & 0,1014 & 0,0938\\
 & \textbf{2} & 0,0014 & 0,0305 & 0,0984 & 0,1762 & 0,2458 & 0,2966 & 0,3241 & 0,3292 & 0,3280 & 0,3110 & 0,2780 & 0,2436 & 0,2344\\
 & \textbf{3} & 0,0000 & 0,0021 & 0,0146 & 0,0415 & 0,0819 & 0,1318 & 0,1852 & 0,2195 & 0,2355 & 0,2765 & 0,3032 & 0,3121 & 0,3125\\
 & \textbf{4} & 0,0000 & 0,0001 & 0,0012 & 0,0055 & 0,0154 & 0,0330 & 0,0595 & 0,0823 & 0,0951 & 0,1382 & 0,1861 & 0,2249 & 0,2344\\
 & \textbf{5} & 0,0000 & 0,0000 & 0,0001 & 0,0004 & 0,0015 & 0,0044 & 0,0102 & 0,0165 & 0,0205 & 0,0369 & 0,0609 & 0,0864 & 0,0938\\
 & \textbf{6} & 0,0000 & 0,0000 & 0,0000 & 0,0000 & 0,0001 & 0,0002 & 0,0007 & 0,0014 & 0,0018 & 0,0041 & 0,0083 & 0,0138 & 0,0156\\
\midrule
\multirow{8}{*}{\textbf{7}} & \textbf{0} & 0,9321 & 0,6983 & 0,4783 & 0,3206 & 0,2097 & 0,1335 & 0,0824 & 0,0585 & 0,0490 & 0,0280 & 0,0152 & 0,0090 & 0,0078\\
 & \textbf{1} & 0,0659 & 0,2573 & 0,3720 & 0,3960 & 0,3670 & 0,3115 & 0,2471 & 0,2048 & 0,1848 & 0,1306 & 0,0872 & 0,0604 & 0,0547\\
 & \textbf{2} & 0,0020 & 0,0406 & 0,1240 & 0,2097 & 0,2753 & 0,3115 & 0,3177 & 0,3073 & 0,2985 & 0,2613 & 0,2140 & 0,1740 & 0,1641\\
 & \textbf{3} & 0,0000 & 0,0036 & 0,0230 & 0,0617 & 0,1147 & 0,1730 & 0,2269 & 0,2561 & 0,2679 & 0,2903 & 0,2918 & 0,2786 & 0,2734\\
 & \textbf{4} & 0,0000 & 0,0002 & 0,0026 & 0,0109 & 0,0287 & 0,0577 & 0,0972 & 0,1280 & 0,1442 & 0,1935 & 0,2388 & 0,2676 & 0,2734\\
 & \textbf{5} & 0,0000 & 0,0000 & 0,0002 & 0,0012 & 0,0043 & 0,0115 & 0,0250 & 0,0384 & 0,0466 & 0,0774 & 0,1172 & 0,1543 & 0,1641\\
 & \textbf{6} & 0,0000 & 0,0000 & 0,0000 & 0,0001 & 0,0004 & 0,0013 & 0,0036 & 0,0064 & 0,0084 & 0,0172 & 0,0320 & 0,0494 & 0,0547\\
 & \textbf{7} & 0,0000 & 0,0000 & 0,0000 & 0,0000 & 0,0000 & 0,0001 & 0,0002 & 0,0005 & 0,0006 & 0,0016 & 0,0037 & 0,0068 & 0,0078\\
\midrule
\multirow{9}{*}{\textbf{8}} & \textbf{0} & 0,9227 & 0,6634 & 0,4305 & 0,2725 & 0,1678 & 0,1001 & 0,0576 & 0,0390 & 0,0319 & 0,0168 & 0,0084 & 0,0046 & 0,0039\\
 & \textbf{1} & 0,0746 & 0,2793 & 0,3826 & 0,3847 & 0,3355 & 0,2670 & 0,1977 & 0,1561 & 0,1373 & 0,0896 & 0,0548 & 0,0352 & 0,0312\\
 & \textbf{2} & 0,0026 & 0,0515 & 0,1488 & 0,2376 & 0,2936 & 0,3115 & 0,2965 & 0,2731 & 0,2587 & 0,2090 & 0,1569 & 0,1183 & 0,1094\\
 & \textbf{3} & 0,0001 & 0,0054 & 0,0331 & 0,0839 & 0,1468 & 0,2076 & 0,2541 & 0,2731 & 0,2786 & 0,2787 & 0,2568 & 0,2273 & 0,2188\\
 & \textbf{4} & 0,0000 & 0,0004 & 0,0046 & 0,0185 & 0,0459 & 0,0865 & 0,1361 & 0,1707 & 0,1875 & 0,2322 & 0,2627 & 0,2730 & 0,2734\\
 & \textbf{5} & 0,0000 & 0,0000 & 0,0004 & 0,0026 & 0,0092 & 0,0231 & 0,0467 & 0,0683 & 0,0808 & 0,1239 & 0,1719 & 0,2098 & 0,2188\\
 & \textbf{6} & 0,0000 & 0,0000 & 0,0000 & 0,0002 & 0,0011 & 0,0038 & 0,0100 & 0,0171 & 0,0217 & 0,0413 & 0,0703 & 0,1008 & 0,1094\\
 & \textbf{7} & 0,0000 & 0,0000 & 0,0000 & 0,0000 & 0,0001 & 0,0004 & 0,0012 & 0,0024 & 0,0033 & 0,0079 & 0,0164 & 0,0277 & 0,0312\\
 & \textbf{8} & 0,0000 & 0,0000 & 0,0000 & 0,0000 & 0,0000 & 0,0000 & 0,0001 & 0,0002 & 0,0002 & 0,0007 & 0,0017 & 0,0033 & 0,0039\\
\midrule
\multirow{10}{*}{\textbf{9}} & \textbf{0} & 0,9135 & 0,6302 & 0,3874 & 0,2316 & 0,1342 & 0,0751 & 0,0404 & 0,0260 & 0,0207 & 0,0101 & 0,0046 & 0,0023 & 0,0020\\
 & \textbf{1} & 0,0830 & 0,2985 & 0,3874 & 0,3679 & 0,3020 & 0,2253 & 0,1556 & 0,1171 & 0,1004 & 0,0605 & 0,0339 & 0,0202 & 0,0176\\
 & \textbf{2} & 0,0034 & 0,0629 & 0,1722 & 0,2597 & 0,3020 & 0,3003 & 0,2668 & 0,2341 & 0,2162 & 0,1612 & 0,1110 & 0,0776 & 0,0703\\
 & \textbf{3} & 0,0001 & 0,0077 & 0,0446 & 0,1069 & 0,1762 & 0,2336 & 0,2668 & 0,2731 & 0,2716 & 0,2508 & 0,2119 & 0,1739 & 0,1641\\
 & \textbf{4} & 0,0000 & 0,0006 & 0,0074 & 0,0283 & 0,0661 & 0,1168 & 0,1715 & 0,2048 & 0,2194 & 0,2508 & 0,2600 & 0,2506 & 0,2461\\
 & \textbf{5} & 0,0000 & 0,0000 & 0,0008 & 0,0050 & 0,0165 & 0,0389 & 0,0735 & 0,1024 & 0,1181 & 0,1672 & 0,2128 & 0,2408 & 0,2461\\
 & \textbf{6} & 0,0000 & 0,0000 & 0,0001 & 0,0006 & 0,0028 & 0,0087 & 0,0210 & 0,0341 & 0,0424 & 0,0743 & 0,1160 & 0,1542 & 0,1641\\
 & \textbf{7} & 0,0000 & 0,0000 & 0,0000 & 0,0000 & 0,0003 & 0,0012 & 0,0039 & 0,0073 & 0,0098 & 0,0212 & 0,0407 & 0,0635 & 0,0703\\
 & \textbf{8} & 0,0000 & 0,0000 & 0,0000 & 0,0000 & 0,0000 & 0,0001 & 0,0004 & 0,0009 & 0,0013 & 0,0035 & 0,0083 & 0,0153 & 0,0176\\
 & \textbf{9} & 0,0000 & 0,0000 & 0,0000 & 0,0000 & 0,0000 & 0,0000 & 0,0000 & 0,0001 & 0,0001 & 0,0003 & 0,0008 & 0,0016 & 0,0020\\
\bottomrule\end{tabular}
\end{center}
\clearpage
\label{tnorm}{\color{fmgreen}\large\bfseries Tabla de la distribución normal $N(0,1)$}\par\smallskip
Da \fmm{F(z)=P(Z\leq z)}, el área a la izquierda de \fmm{z}. La fila da las unidades y las décimas de \fmm{z}; la columna, las centésimas (apartado 10.6).\par\smallskip
\begin{center}\linespread{1}\selectfont\small\setlength{\tabcolsep}{6pt}\renewcommand{\arraystretch}{1.08}
\begin{tabular}{c|cccccccccc}\toprule
$z$ & \textbf{0,00} & \textbf{0,01} & \textbf{0,02} & \textbf{0,03} & \textbf{0,04} & \textbf{0,05} & \textbf{0,06} & \textbf{0,07} & \textbf{0,08} & \textbf{0,09}\\\midrule
\textbf{0,0} & 0,5000 & 0,5040 & 0,5080 & 0,5120 & 0,5160 & 0,5199 & 0,5239 & 0,5279 & 0,5319 & 0,5359\\
\textbf{0,1} & 0,5398 & 0,5438 & 0,5478 & 0,5517 & 0,5557 & 0,5596 & 0,5636 & 0,5675 & 0,5714 & 0,5753\\
\textbf{0,2} & 0,5793 & 0,5832 & 0,5871 & 0,5910 & 0,5948 & 0,5987 & 0,6026 & 0,6064 & 0,6103 & 0,6141\\
\textbf{0,3} & 0,6179 & 0,6217 & 0,6255 & 0,6293 & 0,6331 & 0,6368 & 0,6406 & 0,6443 & 0,6480 & 0,6517\\
\textbf{0,4} & 0,6554 & 0,6591 & 0,6628 & 0,6664 & 0,6700 & 0,6736 & 0,6772 & 0,6808 & 0,6844 & 0,6879\\[2pt]
\textbf{0,5} & 0,6915 & 0,6950 & 0,6985 & 0,7019 & 0,7054 & 0,7088 & 0,7123 & 0,7157 & 0,7190 & 0,7224\\
\textbf{0,6} & 0,7257 & 0,7291 & 0,7324 & 0,7357 & 0,7389 & 0,7422 & 0,7454 & 0,7486 & 0,7517 & 0,7549\\
\textbf{0,7} & 0,7580 & 0,7611 & 0,7642 & 0,7673 & 0,7704 & 0,7734 & 0,7764 & 0,7794 & 0,7823 & 0,7852\\
\textbf{0,8} & 0,7881 & 0,7910 & 0,7939 & 0,7967 & 0,7995 & 0,8023 & 0,8051 & 0,8078 & 0,8106 & 0,8133\\
\textbf{0,9} & 0,8159 & 0,8186 & 0,8212 & 0,8238 & 0,8264 & 0,8289 & 0,8315 & 0,8340 & 0,8365 & 0,8389\\[2pt]
\textbf{1,0} & 0,8413 & 0,8438 & 0,8461 & 0,8485 & 0,8508 & 0,8531 & 0,8554 & 0,8577 & 0,8599 & 0,8621\\
\textbf{1,1} & 0,8643 & 0,8665 & 0,8686 & 0,8708 & 0,8729 & 0,8749 & 0,8770 & 0,8790 & 0,8810 & 0,8830\\
\textbf{1,2} & 0,8849 & 0,8869 & 0,8888 & 0,8907 & 0,8925 & 0,8944 & 0,8962 & 0,8980 & 0,8997 & 0,9015\\
\textbf{1,3} & 0,9032 & 0,9049 & 0,9066 & 0,9082 & 0,9099 & 0,9115 & 0,9131 & 0,9147 & 0,9162 & 0,9177\\
\textbf{1,4} & 0,9192 & 0,9207 & 0,9222 & 0,9236 & 0,9251 & 0,9265 & 0,9279 & 0,9292 & 0,9306 & 0,9319\\[2pt]
\textbf{1,5} & 0,9332 & 0,9345 & 0,9357 & 0,9370 & 0,9382 & 0,9394 & 0,9406 & 0,9418 & 0,9429 & 0,9441\\
\textbf{1,6} & 0,9452 & 0,9463 & 0,9474 & 0,9484 & 0,9495 & 0,9505 & 0,9515 & 0,9525 & 0,9535 & 0,9545\\
\textbf{1,7} & 0,9554 & 0,9564 & 0,9573 & 0,9582 & 0,9591 & 0,9599 & 0,9608 & 0,9616 & 0,9625 & 0,9633\\
\textbf{1,8} & 0,9641 & 0,9649 & 0,9656 & 0,9664 & 0,9671 & 0,9678 & 0,9686 & 0,9693 & 0,9699 & 0,9706\\
\textbf{1,9} & 0,9713 & 0,9719 & 0,9726 & 0,9732 & 0,9738 & 0,9744 & 0,9750 & 0,9756 & 0,9761 & 0,9767\\[2pt]
\textbf{2,0} & 0,9772 & 0,9778 & 0,9783 & 0,9788 & 0,9793 & 0,9798 & 0,9803 & 0,9808 & 0,9812 & 0,9817\\
\textbf{2,1} & 0,9821 & 0,9826 & 0,9830 & 0,9834 & 0,9838 & 0,9842 & 0,9846 & 0,9850 & 0,9854 & 0,9857\\
\textbf{2,2} & 0,9861 & 0,9864 & 0,9868 & 0,9871 & 0,9875 & 0,9878 & 0,9881 & 0,9884 & 0,9887 & 0,9890\\
\textbf{2,3} & 0,9893 & 0,9896 & 0,9898 & 0,9901 & 0,9904 & 0,9906 & 0,9909 & 0,9911 & 0,9913 & 0,9916\\
\textbf{2,4} & 0,9918 & 0,9920 & 0,9922 & 0,9925 & 0,9927 & 0,9929 & 0,9931 & 0,9932 & 0,9934 & 0,9936\\[2pt]
\textbf{2,5} & 0,9938 & 0,9940 & 0,9941 & 0,9943 & 0,9945 & 0,9946 & 0,9948 & 0,9949 & 0,9951 & 0,9952\\
\textbf{2,6} & 0,9953 & 0,9955 & 0,9956 & 0,9957 & 0,9959 & 0,9960 & 0,9961 & 0,9962 & 0,9963 & 0,9964\\
\textbf{2,7} & 0,9965 & 0,9966 & 0,9967 & 0,9968 & 0,9969 & 0,9970 & 0,9971 & 0,9972 & 0,9973 & 0,9974\\
\textbf{2,8} & 0,9974 & 0,9975 & 0,9976 & 0,9977 & 0,9977 & 0,9978 & 0,9979 & 0,9979 & 0,9980 & 0,9981\\
\textbf{2,9} & 0,9981 & 0,9982 & 0,9982 & 0,9983 & 0,9984 & 0,9984 & 0,9985 & 0,9985 & 0,9986 & 0,9986\\[2pt]
\textbf{3,0} & 0,9987 & 0,9987 & 0,9987 & 0,9988 & 0,9988 & 0,9989 & 0,9989 & 0,9989 & 0,9990 & 0,9990\\
\textbf{3,1} & 0,9990 & 0,9991 & 0,9991 & 0,9991 & 0,9992 & 0,9992 & 0,9992 & 0,9992 & 0,9993 & 0,9993\\
\textbf{3,2} & 0,9993 & 0,9993 & 0,9994 & 0,9994 & 0,9994 & 0,9994 & 0,9994 & 0,9995 & 0,9995 & 0,9995\\
\textbf{3,3} & 0,9995 & 0,9995 & 0,9995 & 0,9996 & 0,9996 & 0,9996 & 0,9996 & 0,9996 & 0,9996 & 0,9997\\
\textbf{3,4} & 0,9997 & 0,9997 & 0,9997 & 0,9997 & 0,9997 & 0,9997 & 0,9997 & 0,9997 & 0,9997 & 0,9998\\[2pt]
\textbf{3,5} & 0,9998 & 0,9998 & 0,9998 & 0,9998 & 0,9998 & 0,9998 & 0,9998 & 0,9998 & 0,9998 & 0,9998\\
\textbf{3,6} & 0,9998 & 0,9998 & 0,9999 & 0,9999 & 0,9999 & 0,9999 & 0,9999 & 0,9999 & 0,9999 & 0,9999\\
\textbf{3,7} & 0,9999 & 0,9999 & 0,9999 & 0,9999 & 0,9999 & 0,9999 & 0,9999 & 0,9999 & 0,9999 & 0,9999\\
\textbf{3,8} & 0,9999 & 0,9999 & 0,9999 & 0,9999 & 0,9999 & 0,9999 & 0,9999 & 0,9999 & 0,9999 & 0,9999\\
\textbf{3,9} & 1,0000 & 1,0000 & 1,0000 & 1,0000 & 1,0000 & 1,0000 & 1,0000 & 1,0000 & 1,0000 & 1,0000\\
\bottomrule\end{tabular}
\end{center}

\end{document}
